% Edit the title, individual panel descriptions, and shared notes here.
% All figure placements read these commands; edit once to update every continuation.
\newcommand{\FigSFiveTitle}{Sensitivity of the Scale 3 bottleneck definition and localized frozen-model perturbation effects.}
\newcommand{\FigSFiveA}{Percentage of high-pseudotime 24~h source cells among unique sources of Scale 3 edges retained at the top 0.5\%, 1\%, 2\%, 5\%, and 10\% risk-flux thresholds. The respective unique-source counts are 17, 43, 85, 189, and 351. High pseudotime is defined as pseudotime\_norm at or above the 75th percentile of all 1,043 Walesky 24~h cells. The dashed line marks the 25.0\% background fraction. Percentages use the unique-source count at each flux threshold as the denominator.}
\newcommand{\FigSFiveB}{Mean and maximum absolute TrajRisk-score change among all 910 Walesky 48~h target cells after removing Scale 3 edges from one source-cell pseudotime quartile at a time and rerunning the frozen model. Quartiles Q1--Q4 are defined over all 1,043 cells at 24~h using full-precision, linearly interpolated cutoffs; ties enter the higher quartile (\TableRef{S6F}). Q1--Q4 contain 261, 260, 261, and 261 cells, respectively. The vertical axis is logarithmic; score changes are dimensionless. Mean changes are 0.0015--0.0018, while Q4 yields the largest maximum single-cell change (0.0979). Across all 809 Walesky 96~h cells, mean and maximum absolute changes are zero for every quartile.}
\newcommand{\FigSFiveNotes}{Scale 3 connects 24--48~h. Model weights and node features are retained during each edge-removal forward pass, with no retraining. Source counts, target counts, and maxima describe the model's cell-level outputs rather than independent experimental replicates. The panels report sensitivity summaries and contain no plotted hypothesis tests.}
